% ---------------------------------------------------------------------------
\section{Related work}
\label{sec:related}
% ---------------------------------------------------------------------------

\emph{Neighbourhood overlap.} That adjacent nodes sharing many neighbours tend to share a
community underlies the edge clustering coefficient~\cite{radicchi2004defining}, the
structural similarity of SCAN~\cite{xu2007scan}, shared-near-neighbour
clustering~\cite{jarvis1973clustering} and the link between overlap and tie
strength~\cite{onnela2007structure}; the same indices are the local predictors of link
prediction, where no index is best on every network~\cite{ghasemian2020stacking}.
Jarvis--Patrick clustering, TECTONIC~\cite{tsourakakis2017scalable} and its
refinement~\cite{chen2024parallel} threshold scored edges and take connected components;
our two-test ablation is of this kind, restricted to wedges.

\emph{Clique and triad percolation.} Clique percolation~\cite{palla2005uncovering} joins
$k$-cliques sharing $k-1$ nodes, also on graphs thresholded by link
weight~\cite{kumpula2008sequential,farkas2007weighted}; for $k=3$ these are also the
truss communities~\cite{cohen2008trusses,huang2014querying,sariyuce2015finding}. Our
triangles-only ablation is 3-clique percolation on the graph thresholded at $J\ge t$
(closed-neighbourhood Jaccard). Triad percolation~\cite{zhang2018triad} already admits
open triads, growing seeds until a stopping rule is met; HyperCommon admits a wedge only
when all three of its pairs pass and needs no seeds or stopping rule.

\emph{Link communities.} Ahn, Bagrow and Lehmann~\cite{ahn2010link} call two edges
$\{c,i\}$ and $\{c,j\}$ similar when $N[i]$ and $N[j]$ overlap, with similarity $J(i,j)$,
and cluster edges by single linkage; hyperlink communities extend this to higher-order
networks~\cite{lotito2024hyperlink}. In our language $J(i,j)$ is the commonality of the
third pair of the triple $\{i,c,j\}$, the pair that is not an edge when the wedge is
open. Line and clique graphs~\cite{evans2009line,evans2010clique} frame such
clusterings, and our two-node-overlap graph is the $2$-line graph of the
hypergraph~\cite{aksoy2020hypernetwork}.

\emph{What HyperCommon combines.} It tests the two edges of a wedge, as thresholded
clique percolation and SCAN do, and its non-adjacent pair, as link communities do, and
joins admitted triples through any shared pair. Determinism and nesting in the threshold
are shared with thresholded clique percolation and single-linkage link communities;
what HyperCommon adds is the exact-recovery window (Theorem~\ref{thm:window}), which
ties the threshold to the structure it recovers. To choose the threshold without labels
we use the overlapping modularity of Shen \emph{et al.}~\cite{shen2009detect}, which
peaks near the critical point of $k$-clique percolation~\cite{toth2013overlapping}; for
partitions it is the Newman--Girvan modularity, whose resolution
limit~\cite{fortunato2007resolution} the absolute Potts model~\cite{ronhovde2010local}
and the Constant Potts Model~\cite{traag2011narrow} avoid. The construction was first
formulated, in an order-dependent form, in Ref.~\cite{deinega2023thesis}.
